\begin{center}
\small
\begin{tabular}{@{}lrrrr@{}}
\toprule
Height-one integer quartic & Variables & Terms & Total degree & Individual degree\\
\midrule
Incoming source519 compiler & 7,431 & 34,086 & 4 & 2\\
New matched source519 & 2,988 & 31,218 & 4 & 2\\
New source519, square aliases & 2,776 & 30,554 & 4 & 4\\
New principal on source442 & \textbf{2,215} & \textbf{30,468} & 4 & 4\\
\bottomrule
\end{tabular}
\end{center}
All four rows have coefficient height exactly one and constant coefficient
one. The matched 2,988-variable row reduces the variable count by about
59.8\% and the monomial count by about 8.4\%, with the same source and the same
total and individual degree bounds. The 2,776-variable row reduces the matched
variable count further, allowing fourth powers of individual coordinates.
The 2,215-variable principal result also changes the input to the later
442-variable source; the compiler improvement and source improvement are
therefore not conflated.

The complementary direct arithmetic equations are detailed in
Section~\ref{sec:arithmetic}. Their coefficients are larger but their dimensions
are much smaller. The natural equation with 240 variables has 188 quadratic
residuals, 25,987 collected monomials, degree four, and height 90,589,484.
The one-variable-smaller natural companion has 239 variables, 187 residuals,
360,097 monomials, and height 491,869,488. It is a density tradeoff.
